\documentclass[leqno]{article}
% Shared mark scheme preamble for resources.danieljsmith.org: the pack
% preamble, unchanged, then the scheme's own layer. A scheme's .tex is
% rendered from its YAML (scripts/README.md, Mark schemes), never edited.
\input{\CurrentFilePath/pack_preamble.tex}

% The page: equal side margins, so the table runs the full width of the
% header rule. The footers' right edge follows the text's.
\geometry{left=0.8in, right=0.8in}
\makeatletter
\renewcommand{\djs@footR}[1]{%
  \put(\dimexpr1in+\hoffset+\oddsidemargin+\textwidth\relax,
    -\dimexpr\paperheight-0.32in\relax){\makebox[0pt][r]{#1}}}
\makeatother

\usepackage[export]{adjustbox}
\usepackage{tabularray}
\definecolor{msGrid}{RGB}{176, 172, 164}
\definecolor{msTint}{RGB}{245, 243, 239}

% The header on every page, the front page included, and the booklet label
% under the contents-page title. \djsQuestionPageStyle stays empty.
\newcommand{\djsSchemeHeader}[1]{%
  \djsPackHeader[markscheme]{\textit{Mark scheme} - #1}%
  \renewcommand{\djsPackTitle}{#1}%
  \renewcommand{\djsBookletLabel}{Mark scheme}%
}

% The front page: the contents heading, spaced as tocloft spaces the packs'
% "Questions", then the contents table. A question's name and page number
% both link to its first page.
\newcommand{\msContentsHeading}{%
  \vskip\cftbeforetoctitleskip
  {\noindent\cfttoctitlefont\contentsname\par}%
  \vskip\cftaftertoctitleskip}
\NewTblrEnviron{mscontents}
\SetTblrInner[mscontents]{
  colspec = {Q[wd=26mm, l] X[l] Q[wd=16mm, c] Q[wd=14mm, r]},
  width = \linewidth,
  rowsep = 4pt, colsep = 6pt,
  hlines = {0.4pt, msGrid}, vlines = {0.4pt, msGrid},
  hline{2} = {2.2pt, scarlet},
  row{1} = {font=\small\bfseries},
  row{Z} = {font=\bfseries},
  column{2} = {font=\small},
}
\newcommand{\msContentsQuestion}[1]{\hyperref[q:#1]{Question #1}}
\newcommand{\msContentsPage}[1]{\hyperref[q:#1]{\pageref*{q:#1}}}
\newcommand{\msSep}{\hspace{0.9em}}

% The conventions pages' layout: headings in the documents' own bold, tables
% ruled and padded as the scheme's, roomy lists.
\newcommand{\msHead}[1]{%
  \par\vspace{18pt}%
  {\noindent\bfseries #1\par}%
  \vspace{6pt}}
\NewTblrEnviron{mstblr}
\SetTblrInner[mstblr]{rowsep = 6pt, colsep = 6pt,
  hlines = {0.4pt, msGrid}, vlines = {0.4pt, msGrid}}
\newenvironment{msList}
  {\begin{itemize}[leftmargin=1.2em, itemsep=5pt, topsep=3pt, parsep=0pt]}
  {\end{itemize}}

% A question's first page: the link target and the PDF bookmark, as
% \questionitem gives the packs.
\newcommand{\msQnum}{}
\newcommand{\msQuestion}[1]{%
  \clearpage
  \renewcommand{\msQnum}{#1}%
  \phantomsection\label{q:#1}%
  \addcontentsline{toc}{subsection}{\protect\djsTocEntry{#1}}}

% The scheme's table: Question 16mm, Scheme the rest, Marks 16mm, Notes
% 44mm; every cell ruled in msGrid; the scarlet rule under the head row. A
% long table has no caption or "continued" foot; a continuation page opens
% with a muted line above the repeated head row. The rules across a part's
% rows, the alternative's tint and the total row are the rendered table's.
\NewTblrTheme{ms}{%
  \DefTblrTemplate{firsthead}{default}{}%
  \DefTblrTemplate{middlehead,lasthead}{default}{%
    {\small\itshape\color{djsMuted}Question \msQnum\ continued\par}%
    \vspace{6pt}}%
  \DefTblrTemplate{firstfoot,middlefoot,lastfoot}{default}{}%
}
\NewTblrEnviron{msscheme}
\SetTblrOuter[msscheme]{long, theme=ms, entry=none, label=none}
\SetTblrInner[msscheme]{
  colspec = {Q[wd=16mm, l] X[l] Q[wd=16mm, c] Q[wd=44mm, l]},
  width = \linewidth,
  rowsep = 6pt, colsep = 6pt,
  rowhead = 1,
  vlines = {0.4pt, msGrid},
  hline{1} = {0.4pt, msGrid}, hline{2} = {2.2pt, scarlet},
  hline{Z} = {0.4pt, msGrid},
  column{1} = {font=\bfseries}, column{3} = {font=\bfseries},
  column{4} = {font=\small},
  row{1} = {rowsep=4pt, font=\small\bfseries},
}

% Two marks on a row: each note starts with its code.
\newcommand{\msNote}[2]{\textbf{#1}\enspace #2}
% A full-row figure: the worked solution's own picture, scaled uniformly and
% only down, to the table less the Question column. The picture is drawn
% with the line width it has in the solutions document, at its list depth
% there (the optional argument, 1 by default: a part's working), so one
% sized by \linewidth keeps its size; tabularray sets the cell's \linewidth
% to the cell's measured width, which would shrink it at every pass. The
% widths are measured once, in a solution box at the packs' text width
% (pack_preamble.tex's margins, 0.8in and 1in).
\newlength{\msSpanW}
\makeatletter
\newcommand{\ms@keepLW}[1]{%
  \expandafter\xdef\csname ms@solnLW@#1\endcsname{\the\linewidth}}
\AtBeginDocument{%
  \setlength{\msSpanW}{\dimexpr\textwidth-16mm-28pt\relax}%
  \setbox\z@\vbox{%
    \begin{minipage}{\dimexpr\paperwidth-1.8in\relax}
      \begin{djsSolution}
        \ms@keepLW{0}%
        \begin{enumerate}\item \ms@keepLW{1}%
        \begin{enumerate}\item \ms@keepLW{2}%
        \begin{enumerate}\item \ms@keepLW{3}%
        \end{enumerate}\end{enumerate}\end{enumerate}
      \end{djsSolution}
    \end{minipage}}}
\newcommand{\msFigure}[2][1]{%
  \adjustbox{max width=\msSpanW, max totalheight=0.6\textheight, valign=t}{%
    \hsize=\csname ms@solnLW@#1\endcsname
    \linewidth=\hsize
    #2}}
\makeatother
% A figure mark's feature, in the Scheme column.
\newcommand{\msFeature}[1]{\textit{#1}}
% An alternative's title row.
\newcommand{\msAlternative}[1]{\textit{Alternative:} #1}
% Drafts only: a queried mark's code, the query's why under its note, and a
% flag in its part's place.
\newcommand{\msQueried}[2]{#1~\(\to\) #2?}
\newcommand{\msQuery}[1]{{\itshape\color{djsMuted}Query: #1}}
\newcommand{\msFlag}[2]{\textbf{#1:}\enspace #2}


\djsSchemeHeader{Plane Intersections}
\djsSitePath{Core Pure}{Vectors, Matrices \& Linear Transformations}

\djsLastUpdated{25 October 2026}
\begin{document}
\djsSetQuestionCount{7}
\djsAddMarks{1}{5}
\djsAddMarks{2}{8}
\djsAddMarks{3}{8}
\djsAddMarks{4}{11}
\djsAddMarks{5}{10}
\djsAddMarks{6}{10}
\djsAddMarks{7}{14}
\vspace*{6pt}
\djsContentsTitle
\vspace{30pt}
\msContentsHeading
\begin{mscontents}{}
Question & Marks by part & Marks & Page \\
\msContentsQuestion{1} & (a)\,2\msSep (b)\,3 & 5 & \msContentsPage{1} \\
\msContentsQuestion{2} & (a)\,2\msSep (b)\,3\msSep (c)\,3 & 8 & \msContentsPage{2} \\
\msContentsQuestion{3} & (a)\,6\msSep (b)\,2 & 8 & \msContentsPage{3} \\
\msContentsQuestion{4} & (a)\,2\msSep (b)\,5\msSep (c)(i)\,2\msSep (c)(ii)\,2 & 11 & \msContentsPage{4} \\
\msContentsQuestion{5} & (a)\,4\msSep (b)\,6 & 10 & \msContentsPage{5} \\
\msContentsQuestion{6} & (a)(i)\,2\msSep (a)(ii)\,3\msSep (b)\,5 & 10 & \msContentsPage{6} \\
\msContentsQuestion{7} & (a)\,2\msSep (b)(i)\,6\msSep (b)(ii)\,3\msSep (c)(i)\,1\msSep (c)(ii)\,2 & 14 & \msContentsPage{7} \\
\SetCell[c=2]{r} Total & & 66 & \\
\end{mscontents}
\newpage
% The conventions pages, shared by every pack's mark scheme: one source of
% truth, \input after the front page, with no strand or pack logic. The keys
% on the first side, the rules on the second. The scheme layer defines
% \msHead, msList and mstblr.
\vspace*{6pt}%
{\noindent\cfttoctitlefont Reading the mark scheme\par}%
\vspace{10pt}%
These rules apply throughout, unless a note says otherwise.

\msHead{Marks}
Every mark is worth one, so \textbf{M1 A1} on a row is two marks. Only marks
coded ft follow through.\par
\vspace{10pt}%
\begin{mstblr}{
  colspec = {Q[wd=15mm, l] X[l]},
  width = \linewidth,
  hline{2} = {2.2pt, scarlet},
  column{1} = {font=\bfseries},
  row{1} = {font=\small\bfseries, rowsep=4pt},
}
Code & Meaning \\
M1 & Method: a correct method, applied to this question. A slip doesn't cost it if the method is clear. A formula must be used, not just quoted. \\
dM1 & Dependent method: as M1, but only if the method mark before it was earned, unless its note names another. \\
A1 & Accuracy: a correct answer or step. It needs the method mark it follows. \\
A1ft & As A1, following through from an earlier answer that the note names. \\
A1~AG & The last mark of a part whose answer is printed in the question: complete working, with no errors. \\
B1 & A correct value, fact or statement. It needs no method mark. \\
B1ft & As B1, following through from an earlier answer that the note names. \\
E1 & A written explanation, interpretation or conclusion. Concluding a proof by induction needs every earlier mark. \\
\end{mstblr}

\msHead{Abbreviations}
ft and AG qualify a code in the Marks column. The others appear in the
notes.\par
\vspace{10pt}%
\begin{mstblr}{
  colspec = {Q[wd=17mm, l] X[l] Q[wd=15mm, l] X[l]},
  width = \linewidth,
  column{1} = {font=\bfseries}, column{3} = {font=\bfseries},
}
ft & follow through & AG & answer given in the question \\
cao & correct answer only, with no follow-through & cso & correct solution only: no errors in the part \\
oe & or an equivalent form & awrt & any answer that rounds to this value \\
soi & seen, or implied by later working & isw & ignore working after a correct answer \\
SC & special case: credit outside the main method & their & the student's own earlier value \\
condone & a slip that doesn't cost the mark & [\;\,] & working the mark doesn't need \\
s.f. & significant figures & d.p. & decimal places \\
\end{mstblr}

\newpage
\msHead{Reading the scheme}
\begin{msList}
\item Each row is a line of working, with the marks it earns beside it.
  Where a row carries two marks, each note starts with its code.
\item \textbf{[\textit{n}]} ends each part with its total.
\item A shaded \textit{Alternative} is a second method for the same total.
  Mark an answer by the method it follows.
\item Any other valid method earns the equivalent marks, unless the question
  names a method or says \textit{hence}.
\item A sketch earns one B1 for each feature named below its figure.
\end{msList}

\msHead{Accuracy and answers}
\begin{msList}
\item Unless a question says otherwise, give answers to 3 significant
  figures, and angles in degrees to 1 decimal place. Angles in radians take
  3 significant figures.
\item A final answer with more figures is accepted if it rounds to the
  scheme's value; one with fewer is not. An answer spoiled by rounding too
  early loses its accuracy mark.
\item Where a question asks for an exact answer, a decimal loses the final
  accuracy mark. Otherwise either is accepted.
\item Wrong or missing units cost nothing unless the question asks for them.
\item A correct answer with no working earns full marks, unless the question
  needs the working: show that, prove, hence, a method it names, or an exact
  answer. A wrong answer with no working earns nothing. Calculator answers
  are treated the same way.
\item A correct answer, once reached, isn't spoiled by later working (isw).
\end{msList}

\msHead{Value of $\boldsymbol{g}$}
\begin{msList}
\item Take $g = 9.8\,\mathrm{m\,s^{-2}}$ unless a question states a value,
  which is then the only one accepted. Answers that use $g = 9.8$ may be
  given to 2 or 3 significant figures.
\item Where a question doesn't state $g$, an answer worked with $g = 10$
  throughout earns the same marks, keeping any value printed in the
  question. The worked solutions list the final answers with $g = 10$.
\end{msList}


\msQuestion{1}
\begin{msscheme}{
  hline{3-4,6-8} = {2-Z}{0.4pt, msGrid},
  hline{5,9} = {0.4pt, msGrid},
  row{9} = {font=\bfseries},
}
Question & Scheme & Marks & Notes \\
1(a) & {\(\det\mathbf{A} = 3(-7-0)+2(-14+28)+1(0-7)\)} & M1 & {Expands along any row or column, with three \(2\times2\) minors and their signs. Condone one slip in a minor.} \\
& {\(\phantom{\det\mathbf{A}} = -21+28-7 = 0\)\\[4pt] so \(\mathbf{A}\) is singular and there is no unique solution} & A1~AG & {Correct expansion leading without error to 0, then both conclusions, singular and no unique solution.} \\
\nopagebreak & & [2] & \\
1(b) & {\((3x-2y+z)+2(2x+y-4z) = 5+2(1)\)} & M1 & {Combines two equations to compare with the third, or solves two and substitutes into the third. Condone arithmetic slips.} \\
& {\(\Rightarrow 7x-7z=7\), the third equation,\\[4pt] so the system is consistent} & A1 & {Correct combination matching the third equation, constants included, and consistent stated.} \\
& {The three planes meet in a common line (a sheaf)\\[4pt] \(\left[\mathbf{r} = \begin{pmatrix}1\\-1\\0\end{pmatrix} + \lambda\begin{pmatrix}1\\2\\1\end{pmatrix}\right]\)} & E1 & {States the planes meet in a common line, or form a sheaf. Not for infinitely many solutions alone; the line's equation is not needed.} \\
\nopagebreak & & [3] & \\
\nopagebreak \SetCell[c=2]{r} Total for Question 1 & & 5 & \\
\end{msscheme}

\msQuestion{2}
\begin{msscheme}{
  hline{3-4,6-8,10-12} = {2-Z}{0.4pt, msGrid},
  hline{5,9,13} = {0.4pt, msGrid},
  row{13} = {font=\bfseries},
}
Question & Scheme & Marks & Notes \\
2(a) & {\(\det\mathbf{A} = 1(-7) - 2(-6) - 1(5) = 0\)} & B1 & {\(\det\mathbf{A} = 0\) for the coefficient matrix; a calculator value is accepted.} \\
& {\(\mathbf{A}\) is singular, so no unique solution} & E1 & {States that \(\mathbf{A}\) is singular or not invertible, so there is no unique solution. Needs \(\det\mathbf{A} = 0\).} \\
\nopagebreak & & [2] & \\
2(b) & {\((x+2y-z) + (2x-y+4z) = 3 + 1\)} & M1 & {Adds the first two equations, or eliminates one variable from the system by any method, and compares with the third.} \\
& {\(\Rightarrow 3x+y+3z = 4\), the third equation, so consistent} & A1 & {Correct sum identified with the third equation, or \(0 = 0\) by elimination, and a conclusion of consistency.} \\
& {The three planes meet in a common line (a sheaf)} & E1 & {States the planes meet in a single common line, or form a sheaf. Not for infinitely many solutions alone.} \\
\nopagebreak & & [3] & \\
2(c) & {First two equations \(\Rightarrow 3x+y+3z = 4\)} & M1 & {Combines the first two equations to give \(3x+y+3z = 4\), or eliminates by any method. May be seen in (b).} \\
& {\(3x+y+3z = a \ne 4\): contradiction, so inconsistent} & A1 & {Contradiction with the third equation for \(a \ne 4\) and a conclusion of inconsistency.} \\
& {Normals \((1, 2, -1)\), \((2, -1, 4)\), \((3, 1, 3)\): no two parallel\\[4pt] \(\Rightarrow\) triangular prism} & E1 & {States triangular prism, with a reason no two planes are parallel, such as no two normals being multiples.} \\
\nopagebreak & & [3] & \\
\nopagebreak \SetCell[c=2]{r} Total for Question 2 & & 8 & \\
\end{msscheme}

\msQuestion{3}
\begin{msscheme}{
  hline{3-8,10-11} = {2-Z}{0.4pt, msGrid},
  hline{9,12} = {0.4pt, msGrid},
  row{12} = {font=\bfseries},
}
Question & Scheme & Marks & Notes \\
3(a) & {\(\begin{vmatrix}1&-1&2\\2&1&1\\3&0&m+1\end{vmatrix}\)\\[4pt] \(\phantom{A} = 1(m+1) + 1\bigl(2(m+1) - 3\bigr) + 2(0 - 3)\)} & M1 & {Expands the determinant of the coefficient matrix in terms of \(m\) by any row or column. Condone sign slips.} \\
& {\(\phantom{A} = 3(m - 2)\)} & A1 & {\(3(m-2)\) oe, such as \(3m - 6\).} \\
& {\(m \ne 2 \Rightarrow \det \ne 0\): the planes meet in a single point} & E1 & {Relates the non-zero determinant to a unique point of intersection. The point need not be found.} \\
& {\(m = 2\): first \(+\) second gives \(3x + 3z = 7\)\\[4pt] third: \(3x + 3z = 7\)} & M1 & {Substitutes \(m = 2\) and tests consistency, by combining the first two equations to compare with the third, or by elimination.} \\
& {Third equation is the sum of the first two, so consistent} & A1 & {Third equation shown to be the sum of the first two, or elimination reaching \(0 = 0\), with consistency stated or implied.} \\
& {The planes meet in a common line (a sheaf)} & E1 & {States the planes share one common line, or form a sheaf, from a correct consistency result. The line need not be found.} \\
\nopagebreak & & [6] & \\
3(b) & {\(\begin{vmatrix}1&-1&2\\2&1&1\\3&0&4\end{vmatrix} = 1(4) + 1(8 - 3) + 2(0 - 3)\)} & M1 & {Forms the determinant or scalar triple product of the three vectors and attempts to evaluate it, or uses another valid coplanarity test. Condone sign slips.} \\
& {\(= 3 \ne 0\), so not coplanar: the student is incorrect} & A1 & {Value \(\pm 3\), depending on order, and concludes the student is incorrect.} \\
\nopagebreak & & [2] & \\
\nopagebreak \SetCell[c=2]{r} Total for Question 3 & & 8 & \\
\end{msscheme}

\msQuestion{4}
\begin{msscheme}{
  hline{3-4,6-16,18-19,21-22} = {2-Z}{0.4pt, msGrid},
  hline{5,17,20,23} = {0.4pt, msGrid},
  cell{10-15}{2-4} = {bg=msTint},
  row{23} = {font=\bfseries},
}
Question & Scheme & Marks & Notes \\
4(a) & {\(\det\mathbf{M} = 1(k + 1) - 1(3 - 2) + 0\)} & M1 & {Expands the \(3\times3\) determinant along any row or column, with every nonzero \(2\times2\) minor attempted. Condone sign slips.} \\
& {\(\det\mathbf{M} = k\), so an inverse exists for \(k \neq 0\)} & A1 & \\
\nopagebreak & & [2] & \\
4(b) & {\(\mathbf{M}^{-1} = \frac12\begin{pmatrix}3 & -1 & 1\\-1 & 1 & -1\\-7 & 3 & -1\end{pmatrix}\)} & B1 & {Correct inverse with \(k = 2\), from a calculator or otherwise.} \\
& {\(\begin{pmatrix}x\\y\\z\end{pmatrix} = \their{\mathbf{M}^{-1}}\begin{pmatrix}1\\p\\4\end{pmatrix} = \ldots\)} & M1 & {Multiplies their \(3\times3\) inverse by \(\begin{pmatrix}1\\p\\4\end{pmatrix}\), in that order, for at least one component.} \\
& {\(x = \frac{7-p}{2}\)} & A1 & {oe. Coordinates may be given as a column vector.} \\
& {\(y = \frac{p-5}{2}\)} & A1 & \\
& {\(z = \frac{3p-11}{2}\)} & A1 & \\
& \SetCell[c=3]{l} \msAlternative{Elimination} & & \\
\nopagebreak & {\(y = 1 - x \Rightarrow 3x + 2(1 - x) + z = p\)\\[4pt] \(\phantom{y = 1 - x \Rightarrow}\ 2x - (1 - x) + z = 4\)} & M1 & {Eliminates one variable from two pairs of the equations, giving two equations in the other two. Condone sign slips.} \\
& {\(x + z = p - 2,\quad 3x + z = 5\)} & A1 & {oe, or the correct pair for another eliminated variable.} \\
& {\(2x = 5 - (p - 2) \Rightarrow x = \ldots\)\\[4pt] \(y = 1 - x = \ldots,\quad z = 5 - 3x = \ldots\)} & dM1 & {Solves their pair for one variable and substitutes back to find the other two.} \\
& {\(x = \frac{7-p}{2}\)} & A1 & {oe. Coordinates may be given as a column vector.} \\
& {\(y = \frac{p-5}{2},\quad z = \frac{3p-11}{2}\)} & A1 & {Both correct, oe.} \\
\nopagebreak & & [5] & \\
4(c)(i) & {\((3x + z) - (x + y) = 2x - y + z \Rightarrow q - 1 = 4\)} & M1 & {Combines two of the equations to match the variable terms of the third, or eliminates variables to compare with it, reaching an equation in \(q\). Condone sign slips.} \\
& {\(q = 5\)} & A1 & \\
\nopagebreak & & [2] & \\
4(c)(ii) & {Third equation follows from the first two, so infinitely many solutions} & E1 & {States that one equation is a combination of the other two, or that the system is consistent with no unique solution, and so has infinitely many solutions.} \\
& {The three planes meet in a common line (a sheaf)} & E1 & {States that the planes meet in a single common line; ``sheaf'' alone accepted. The line's equation is not needed.} \\
\nopagebreak & & [2] & \\
\nopagebreak \SetCell[c=2]{r} Total for Question 4 & & 11 & \\
\end{msscheme}

\msQuestion{5}
\begin{msscheme}{
  hline{3-6,8-13} = {2-Z}{0.4pt, msGrid},
  hline{7,14} = {0.4pt, msGrid},
  row{14} = {font=\bfseries},
}
Question & Scheme & Marks & Notes \\
5(a) & {\(\begin{vmatrix}1 & -1 & 2\\ k+1 & k & 2-k\\ 5-2k & 2-2k & k+2\end{vmatrix} = 0\)\\[4pt] \(\Rightarrow k(k+2)-(2-k)(2-2k) + \ldots = 0\)} & M1 & {Sets the determinant of the coefficient matrix to zero and expands it, with at least two of the three \(2 \times 2\) terms correct in form. Condone sign slips.} \\
& {\(-2k^2 + 10k - 8 = 0\)} & A1 & {Correct quadratic, oe such as \(k^2 - 5k + 4 = 0\).} \\
& {\(-2(k-1)(k-4) = 0 \Rightarrow k = \ldots\)} & dM1 & {Solves their three-term quadratic by any means, a calculator allowed.} \\
& {\(k = 1\) or \(k = 4\)} & A1 & {Both values and no others.} \\
\nopagebreak & & [4] & \\
5(b) & {\(k = 1\): \((x - y + 2z) + (2x + y + z) = 4 + 3\)\\[4pt] \(\phantom{k = 1\text{:}} \Rightarrow 3x + 3z = 7\)} & M1 & {Substitutes \(k = 1\) and combines two equations to compare with the third, or eliminates a variable from two pairs and compares.} \\
& {Same as the third equation, so consistent} & A1 & {Shows the third equation is the sum of the first two and states consistent.} \\
& {\(\mathbf{n}_1 = \begin{pmatrix}1\\-1\\2\end{pmatrix}\), \(\mathbf{n}_2 = \begin{pmatrix}2\\1\\1\end{pmatrix}\) not parallel\\[4pt] \(\Rightarrow\) the planes meet in a common line: a sheaf} & E1 & {States sheaf, or a common line, with at least two normals shown not parallel.} \\
& {\(k = 4\): \(2(x - y + 2z) - (5x + 4y - 2z) = 2 \times 4 - 0\)\\[4pt] \(\phantom{k = 4\text{:}} \Rightarrow -3x - 6y + 6z = 8\)} & M1 & {Substitutes \(k = 4\) and combines two equations to compare with the third, or eliminates a variable from two pairs and compares. Condone arithmetic slips.} \\
& {But \(-3x - 6y + 6z = 4\): \(8 \neq 4\), so inconsistent} & A1 & {Correct contradiction and states inconsistent.} \\
& {\(\begin{pmatrix}1\\-1\\2\end{pmatrix}\), \(\begin{pmatrix}5\\4\\-2\end{pmatrix}\), \(\begin{pmatrix}-3\\-6\\6\end{pmatrix}\): no two parallel\\[4pt] \(\Rightarrow\) triangular prism} & E1 & {States triangular prism, with all three pairs of normals shown or stated not parallel. E0 without the normals.} \\
\nopagebreak & & [6] & \\
\nopagebreak \SetCell[c=2]{r} Total for Question 5 & & 10 & \\
\end{msscheme}

\msQuestion{6}
\begin{msscheme}{
  hline{3-4,6-8,10-14} = {2-Z}{0.4pt, msGrid},
  hline{5,9,15} = {0.4pt, msGrid},
  row{15} = {font=\bfseries},
}
Question & Scheme & Marks & Notes \\
6(a)(i) & {\(\det\mathbf{A} = 1(2-b) - a(4-3) = 0\)} & M1 & {Sets an expansion of the determinant of the coefficient matrix equal to zero. Condone one sign slip.} \\
& {\(b = 2-a\)} & A1 & {oe} \\
\nopagebreak & & [2] & \\
6(a)(ii) & {\(x = -1-ay,\quad z = 5+(2a-1)y\)} & M1 & {Finds the line of intersection of two planes, or a point on it, or reduces to two equations in two variables, or writes one plane as a combination of the others. Condone sign slips.} \\
& {\(3(-1-ay)+(2-a)y+2\bigl(5+(2a-1)y\bigr) = c\)} & dM1 & {Uses their \(b\) in terms of \(a\) with their line, point, reduced equations or combination to find \(c\).} \\
& {\(c = 7\)} & A1 & \\
\nopagebreak & & [3] & \\
6(b) & {\(x = -1-ay \Rightarrow 2(-1-ay)+y+z = 3\)} & M1 & {Eliminates a variable between two of the equations, giving one variable in terms of the others. Condone sign slips.} \\
& {\(z = 5+(2a-1)y\)} & A1 & {oe, or the correct result of another elimination.} \\
& {\(3(-1-ay)+(1-a)y+2\bigl(5+(2a-1)y\bigr) = 5\)\\[4pt] \(\Rightarrow 7-y = 5 \Rightarrow y = \ldots\)} & dM1 & {Substitutes their expressions into the third equation, with \(b = 1-a\) and \(c = 5\), and solves for the remaining variable.} \\
& {\(y = 2\)} & A1 & \\
& {\(x = -1-2a,\quad z = 4a+3\)} & A1 & {Both, oe, as coordinates or separate values.} \\
\nopagebreak & & [5] & \\
\nopagebreak \SetCell[c=2]{r} Total for Question 6 & & 10 & \\
\end{msscheme}

\msQuestion{7}
\begin{msscheme}{
  hline{3-4,6-10,12-14,16,18-19} = {2-Z}{0.4pt, msGrid},
  hline{5,11,15,17,20} = {0.4pt, msGrid},
  row{20} = {font=\bfseries},
}
Question & Scheme & Marks & Notes \\
7(a) & {\(\det\mathbf{A} = 3(3 - 0) - 2(1 - 0) + a(-1 - 6)\)} & M1 & {Expands along any row or column, three terms each an element times its \(2 \times 2\) determinant, or uses Sarrus' rule. Condone sign slips.} \\
& {\(\det\mathbf{A} = 7(1 - a)\)} & A1 & {oe, such as \(7 - 7a\).} \\
\nopagebreak & & [2] & \\
7(b)(i) & {Minors: \(\begin{pmatrix}3 & 1 & -7\\ 2+a & 3-2a & -7\\ -3a & -a & 7\end{pmatrix}\)} & M1 & {Attempts all nine \(2 \times 2\) minors; implied by a matrix of cofactors.} \\
& {\(\mathrm{adj}\,\mathbf{A} = \begin{pmatrix}3 & -a-2 & -3a\\ -1 & 3-2a & a\\ -7 & 7 & 7\end{pmatrix}\)} & M1 & {Applies the alternating signs and transposes their matrix of minors, in either order.} \\
& {\(\mathbf{A}^{-1} = \frac{1}{7(1-a)}\begin{pmatrix}3 & -a-2 & -3a\\ -1 & 3-2a & a\\ -7 & 7 & 7\end{pmatrix}\)} & M1 A1 & {\msNote{M1}{Divides their adjugate by their determinant from (a), in terms of \(a\).}\\ \msNote{A1}{Correct inverse, oe; the factor may be multiplied in.}} \\
& {\(\begin{pmatrix}x\\ y\\ z\end{pmatrix} = \their{\mathbf{A}^{-1}}\begin{pmatrix}8\\ 5\\ 3\end{pmatrix}\)\\[4pt] \(\phantom{\begin{pmatrix}x\\ y\\ z\end{pmatrix}} = \frac{1}{7(1-a)}\begin{pmatrix}14-14a\\ 7-7a\\ 0\end{pmatrix}\)} & M1 & {Multiplies their \(\mathbf{A}^{-1}\), in terms of \(a\), by \((8, 5, 3)^{T}\) in that order, at least one element attempted.} \\
& {\(x = 2,\ y = 1,\ z = 0\)} & A1 & {cso; all three values.} \\
\nopagebreak & & [6] & \\
7(b)(ii) & {With \(a = 1\), adds the second and third equations:\\[4pt] \((x + 3y) + (2x - y + z) = 5 + 3\)} & M1 & {Sets \(a = 1\) and compares a combination of two equations with the third.} \\
& {\(\Rightarrow 3x + 2y + z = 8\), the first equation, so consistent} & A1 & {Shows the first equation is the sum of the others, oe, and concludes the equations are consistent or have solutions.} \\
& {The three planes meet in a common line: a sheaf} & E1 & {States the planes share a single common line, or form a sheaf. No check that the normals are non-parallel is needed.} \\
\nopagebreak & & [3] & \\
7(c)(i) & {\(a > 1\)} & B1 & \\
\nopagebreak & & [1] & \\
7(c)(ii) & {\(|7(1-a)|(a^2 + 3) < a^2 + 3 \Rightarrow |7(1-a)| < 1\)} & M1 & {Uses the modulus of their determinant from (a) as the volume scale factor, reaching \(|\det\mathbf{A}| < 1\) or both cases. Not for \(7(1-a) < 1\) alone.} \\
& {\(\frac67 < a < \frac87\)} & A1 & {oe, such as \(\left(\frac67, \frac87\right)\).} \\
\nopagebreak & & [2] & \\
\nopagebreak \SetCell[c=2]{r} Total for Question 7 & & 14 & \\
\end{msscheme}

\end{document}
